\documentclass[11pt,a4paper]{article}

% ==============================================================================
% EXPERIMENT 04: PRINTABLE SOURCE CODE & OUTPUT (EXACT SINGLE 1-PAGE DOCUMENT)
% Formatted with 1.0-inch lateral margins, large readable fonts, and tight top/bottom margins.
% ==============================================================================
\usepackage[utf8]{inputenc}
\usepackage[T1]{fontenc}
\usepackage{lmodern}
\usepackage[left=1.0in, right=1.0in, top=0.32in, bottom=0.32in]{geometry}
\usepackage{amsmath,amssymb}
\usepackage{listings}
\usepackage{tcolorbox}
\tcbuselibrary{skins}
\usepackage{booktabs}
\usepackage{tabularx}
\usepackage{titlesec}
\usepackage{caption}
\usepackage{microtype}
\usepackage{float}

% Disable all running headers, footers, and page numbers for clean print & paste
\pagestyle{empty}
\setlength{\parindent}{0pt}

% Section Titles Formatting (Enlarged)
\titleformat{\section}
  {\normalfont\large\bfseries\raggedright}{\thesection}{0.6em}{}[\titlerule]

\titlespacing*{\section}{0pt}{4pt}{2pt}

% Code Listing Styling (Strict Monochrome with Enlarged Font for High-Contrast Print)
\lstdefinestyle{bwArduino}{
    language=C++,
    basicstyle=\ttfamily\footnotesize,
    keywordstyle=\bfseries,
    commentstyle=\itshape,
    stringstyle=\ttfamily,
    numbers=left,
    numberstyle=\tiny,
    numbersep=6pt,
    stepnumber=1,
    breaklines=true,
    breakatwhitespace=true,
    frame=single,
    rulecolor=\color{black},
    tabsize=2,
    showstringspaces=false,
    columns=fullflexible,
    keepspaces=true,
    lineskip=-0.1pt,
    aboveskip=2pt,
    belowskip=2pt
}

\lstdefinestyle{bwTerminal}{
    basicstyle=\ttfamily\footnotesize,
    numbers=none,
    breaklines=true,
    breakatwhitespace=true,
    frame=single,
    rulecolor=\color{black},
    columns=fullflexible,
    keepspaces=true,
    lineskip=1.2pt,
    aboveskip=2pt,
    belowskip=2pt
}

\captionsetup{font=small,labelfont=bf,skip=2pt}

\begin{document}

% ==============================================================================
% DOCUMENT HEADER
% ==============================================================================
\begin{center}
    {\Large\bfseries Experiment 04: Interfacing HC-SR04 Ultrasonic Sensor with Arduino UNO}\\[0.10cm]
    {\normalsize\textbf{Non-Contact Distance Measurement --- Program Source Code \& Experimental Output}}
\end{center}
\vspace{-0.20cm}
\rule{\textwidth}{0.9pt}
\vspace{0.06cm}

% ==============================================================================
% 1. PROGRAM SOURCE CODE
% ==============================================================================
\section*{1. Arduino C++ Source Code}

\begin{lstlisting}[style=bwArduino, title={\textbf{Listing 4.1: Microcontroller Firmware for Ultrasonic Distance Measurement.}}]
/*
 * Experiment 04: Ultrasonic Distance Measurement using HC-SR04 and Arduino UNO
 * Hardware Pin Connections:
 *   - Arduino Digital Pin 8 -> HC-SR04 TRIG (Trigger Output)
 *   - Arduino Digital Pin 7 -> HC-SR04 ECHO (Echo Pulse Input)
 *   - Arduino 5V & GND     -> HC-SR04 VCC & GND Power Supply
 * Serial Baud Rate: 9600 bps
 */
const int trigPin = 8;    // Pin D8 connected to HC-SR04 TRIG terminal
const int echoPin = 7;    // Pin D7 connected to HC-SR04 ECHO terminal
long duration;            // Holds raw Time-of-Flight (ToF) in microseconds
int distance;             // Holds calculated one-way distance in centimeters

void setup() {
  pinMode(trigPin, OUTPUT);  // Configure Trigger pin as digital output
  pinMode(echoPin, INPUT);   // Configure Echo pin as digital input
  Serial.begin(9600);        // Initialize UART Serial Interface at 9600 Baud
}

void loop() {
  // Step 1: Ensure a clean LOW signal before generating trigger pulse
  digitalWrite(trigPin, LOW);
  delayMicroseconds(2);
  // Step 2: Transmit 10 microsecond Active-HIGH Trigger Pulse
  digitalWrite(trigPin, HIGH);
  delayMicroseconds(10);
  digitalWrite(trigPin, LOW);
  // Step 3: Measure HIGH pulse duration on ECHO pin (Time-of-Flight in us)
  duration = pulseIn(echoPin, HIGH);
  // Step 4: Calculate one-way distance in centimeters (Speed of sound = 0.034 cm/us)
  distance = duration * 0.034 / 2;
  // Step 5: Format and stream telemetry data to Serial Monitor
  Serial.print("distance : ");
  Serial.println(distance);
  // Step 6: Measurement cycle delay to prevent acoustic multipath echo interference
  delay(100);
}
\end{lstlisting}

\vspace{0.06cm}

% ==============================================================================
% 2. EXPERIMENTAL OUTPUT & SERIAL TELEMETRY (SERIAL: 2.A -> 2.B)
% ==============================================================================
\section*{2. Experimental Output \& Serial Monitor Telemetry}

\noindent
\textbf{A. Ultrasonic Distance Measurement \& Calibration Log:}\par\vspace{0.06cm}
{\footnotesize
\setlength{\tabcolsep}{6pt}
\renewcommand{\arraystretch}{1.18}
\begin{tabularx}{\textwidth}{|c|c|c|X|}
\hline
\textbf{Trial} & \textbf{Actual Range} & \textbf{Serial Output} & \textbf{Measurement State \& Hardware Observation} \\
\hline
1 & $5.0\,\text{cm}$ & \texttt{distance : 5} & Exact Setpoint Match ($100\%$ Correlation) \\
2 & $5.0\,\text{cm}$ & \texttt{distance : 4} & Stable Reading with Tolerable Noise ($\Delta = -1.0\,\text{cm}$) \\
3 & $5.0\,\text{cm}$ & \texttt{distance : 5} & Exact Setpoint Match ($100\%$ Correlation) \\
4 & $10.0\,\text{cm}$ & \texttt{distance : 10} & Linear Scaling Validated ($100\%$ Correlation) \\
\hline
\end{tabularx}
}

\vspace{0.22cm}

\noindent
\textbf{B. Serial Monitor Stream (9600 Baud):}\par\vspace{0.05cm}
\begin{lstlisting}[style=bwTerminal]
distance : 5
distance : 4
distance : 5
distance : 10
\end{lstlisting}

\end{document}
